
\medskip
\phantomsection\label{prop:sga3-VIB-11.10.bis}%
\noindent\textbf{Proposition 11.10.bis.---}%
\footnote{\label{note:sga3-VIB-119}Editor's note: this
proposition has been added; it is a special case of
\textup{[Th87]}, \SGARefLink{prop:sga3-VIB-1.4}{1.4}--\SGARefLink{lem:sga3-VIB-1.5}{1.5}. The author there refers to an
argument of Deligne \textup{(cf.\ \textup{[Kn71]},
III, Th.\ \SGARefLink{par:sga3-VIB-1.1}{1.1})}; one may also note its similarity with
Serre's argument \textup{(\textup{[Se68]}, Prop.\ 2)}
recalled in \hyperref[prop:sga3-VIB-11.8-bis]{11.8.bis}.}
\emph{Let \(S\) be a Noetherian scheme, let \(G\) be an
\(S\)-group that is flat, quasi-compact, and quasi-separated
over \(S\), let \(\mathcal E\) be a quasi-coherent
\(G\)-\(\mathcal O_S\)-module, and let \(\mathcal M\) be a
coherent \(\mathcal O_S\)-submodule of \(\mathcal E\). Then
\(\mathcal M\) is contained in a coherent
\(\mathcal O_S\)-submodule stable under \(G\).}

\emph{Proof.}
Write \(f\) for the morphism \(G\to S\), and \(\tau\) for
the adjunction morphism
\[
  \mathcal E\longrightarrow f_*f^*(\mathcal E).
\]
By \hyperref[par:sga3-VIB-11.6]{11.6}, the
\(G\)-\(\mathcal O_S\)-module structure on \(\mathcal E\)
is given by an automorphism \(\theta\) of the
\(\mathcal O_G\)-module \(f^*(\mathcal E)\), such that the
morphism
\[
  \delta=\theta\circ\tau
\]
satisfies conditions
\hyperref[diag:sga3-VIB-11.6-1]{\textup{(1)}} and
\hyperref[diag:sga3-VIB-11.6-2]{\textup{(2)}} of
\hyperref[par:sga3-VIB-11.6]{11.6}. (The situation considered
in \textup{[Th87]} is more general: the author considers a
\(G\)-scheme \(X\) and a \(G\)-equivariant
\(\mathcal O_X\)-module \(\mathcal E\)
\textup{(cf.\ \SGARefLink{sga3:i:section:6}{Expos\'e I, Section 6})}; here \(X=S\), endowed
with the trivial action of \(G\).)

Since \(S\) is Noetherian, \(\mathcal E\) is the filtered
direct limit of its coherent submodules
\(\mathcal F_\alpha\) \textup{(cf.\ EGA~I, 9.4.9)}. Then
\(f^*(\mathcal E)\) is the filtered direct limit of the
\(f^*(\mathcal F_\alpha)\), which are submodules of
\(f^*(\mathcal E)\) because \(f\) is flat. Moreover, since
\(f\) is quasi-compact and quasi-separated, by
\hyperref[par:sga3-VIB-11.0]{11.0} the
\(\mathcal O_S\)-module \(f_*f^*(\mathcal E)\) is
quasi-coherent and is the filtered direct limit of the
quasi-coherent submodules \(f_*f^*(\mathcal F_\alpha)\).
Consequently, \(\mathcal E\) is the filtered direct limit of
the quasi-coherent \(\mathcal O_S\)-submodules
\(\mathcal E_\alpha\), where \(\mathcal E_\alpha\) denotes
the inverse image under \(\delta\) of
\(f_*f^*(\mathcal F_\alpha)\); equivalently, it is the kernel
of the composite
\phantomsection\label{form:sga3-VIB-11.10bis-delta-alpha}%
\[
  \overline{\delta}_\alpha:\quad
  \mathcal E\xrightarrow{\ \delta\ }f_*f^*(\mathcal E)
    \longrightarrow f_*f^*(\mathcal E/\mathcal F_\alpha).
\]
Since \(\mathcal M\) is coherent and \(S\) is Noetherian,
every increasing sequence of submodules of \(\mathcal M\)
is stationary; hence \(\mathcal M\) is contained in some
\(\mathcal E_\alpha\). We shall show that every
\(\mathcal E_\alpha\) is coherent and \(G\)-stable.

Let
\[
  u:f^*(\mathcal E)\longrightarrow\varepsilon_*(\mathcal E)
\]
be the morphism corresponding by adjunction to the identity
morphism from \(\mathcal E\) to
\(f_*\varepsilon_*(\mathcal E)=\mathcal E\). Then the
composite
\phantomsection\label{form:sga3-VIB-11.10bis-adjunction-retract}%
\[
  \mathcal E\xrightarrow{\ \delta\ }f_*f^*(\mathcal E)
    \xrightarrow{\ f_*(u)\ }f_*\varepsilon_*(\mathcal E)
    =\mathcal E
\]
is the identity%
\footnote{\label{source-note:sga3-VIB-11.10bis-delta-arrows}Translator's note: the French text labels the first arrow in this composite and the lower
horizontal arrow in the following diagram by \(\tau\). They must be
\(\delta=\theta\circ\tau\): condition 11.6\textup{(2)} gives
\(f_*(u)\circ\delta=\operatorname{id}_{\mathcal E}\), and the definition
of \(\mathcal E_\alpha\) makes the square commute.}
\textup{(cf.\ \hyperref[diag:sga3-VIB-11.6-2]{11.6\textup{(2)}})}.
We have the commutative diagram
\begin{center}
  \phantomsection\label{diag:sga3-VIB-11.10bis-coherence}%
  \begin{tikzcd}[row sep=large,column sep=large]
    \mathcal E_\alpha
      \arrow[r,"\delta"]
      \arrow[d,"i_\alpha"'] &
    f_*f^*(\mathcal F_\alpha)
      \arrow[r,"f_*(u)"]
      \arrow[d] &
    \mathcal F_\alpha
      \arrow[d,"j_\alpha"] \\
    \mathcal E
      \arrow[r,"\delta"'] &
    f_*f^*(\mathcal E)
      \arrow[r,"f_*(u)"'] &
    \mathcal E\mathpunct{,}
  \end{tikzcd}
\end{center}
where \(i_\alpha\) \textup{(respectively \(j_\alpha\))}
denotes the inclusion of \(\mathcal E_\alpha\)
\textup{(respectively \(\mathcal F_\alpha\))} in
\(\mathcal E\). It follows that \(i_\alpha\) factors through
\(j_\alpha\); thus \(\mathcal E_\alpha\) is a submodule of
\(\mathcal F_\alpha\), and is therefore coherent.

Finally, let us show that \(\mathcal E_\alpha\) is
\(G\)-stable
\textup{(cf.\ \hyperref[rem:sga3-VIB-11.7.bis]{11.7.bis})}.
Let \(h:G\to S\) denote a second copy of \(f:G\to S\), and
consider the following commutative diagram:
\begin{center}
  \phantomsection\label{diag:sga3-VIB-11.10bis-projections}%
  \begin{tikzcd}[row sep=large,column sep=large]
    G\times_S G
      \arrow[r,"q"]
      \arrow[d,"p"']
      \arrow[dr,"\varphi" description] &
    G \arrow[d,"h"] \\
    G \arrow[r,"f"'] &
    S\mathpunct{.}
  \end{tikzcd}
\end{center}
The exact sequence
\phantomsection\label{seq:sga3-VIB-11.10bis-first-exact}%
\[
  0\longrightarrow\mathcal E_\alpha
    \longrightarrow\mathcal E
    \xrightarrow{\ \overline{\delta}_\alpha\ }
    h_*h^*(\mathcal E/\mathcal F_\alpha)
\]
therefore gives, since \(f\) is flat, the exact sequence
\phantomsection\label{seq:sga3-VIB-11.10bis-dagger}%
\[
  \tag{\(\dagger\)}
  0\longrightarrow f_*f^*(\mathcal E_\alpha)
    \longrightarrow f_*f^*(\mathcal E)
    \xrightarrow{\ f_*f^*(\overline{\delta}_\alpha)\ }
    f_*f^*h_*h^*(\mathcal E/\mathcal F_\alpha).
\]
Moreover, since \(h\) is quasi-compact and quasi-separated
and \(f\) is flat, the canonical morphism
\phantomsection\label{form:sga3-VIB-11.10bis-flat-base-change}%
\[
  f^*h_*h^*(\mathcal E/\mathcal F_\alpha)
    \longrightarrow
  p_*q^*h^*(\mathcal E/\mathcal F_\alpha)
    =p_*\varphi^*(\mathcal E/\mathcal F_\alpha)
\]
is an isomorphism. Thus the right-hand term in
\hyperref[seq:sga3-VIB-11.10bis-dagger]{\textup{(\(\dagger\))}}
is \(\varphi_*\varphi^*(\mathcal E/\mathcal F_\alpha)\).
Using the notation of \hyperref[par:sga3-VIB-11.6]{11.6}
and writing \(\pi_\alpha\) for the projection
\(\mathcal E\to\mathcal E/\mathcal F_\alpha\), we obtain the
following commutative diagram, whose bottom row is exact:
\begin{center}
  \phantomsection\label{diag:sga3-VIB-11.10bis-stability}%
  {\small
  \begin{tikzcd}[row sep=large,column sep=large]
    {} &
    \mathcal E_\alpha
      \arrow[r,"\delta"]
      \arrow[d,"\delta"'] &
    f_*f^*(\mathcal E)
      \arrow[d,"{f_*(f^*(\pi_\alpha)\boxtimes\delta_G)}"] \\
    f_*f^*(\mathcal E_\alpha)
      \arrow[r] &
    f_*f^*(\mathcal E)
      \arrow[r,"{f_*f^*(\overline{\delta}_\alpha)}"'] &
    \varphi_*\varphi^*(\mathcal E/\mathcal F_\alpha)
  \end{tikzcd}}
\end{center}
Since
\[
  f_*\bigl(f^*(\pi_\alpha)\boxtimes\delta_G\bigr)\circ\delta
\]
vanishes on \(\mathcal E_\alpha\), it follows that
\(\delta\) carries \(\mathcal E_\alpha\) into
\(f_*f^*(\mathcal E_\alpha)\); in other words,
\(\mathcal E_\alpha\) is \(G\)-stable. This proves the
proposition.

\medskip
\phantomsection\label{rem:sga3-VIB-11.10.1}%
\noindent\textbf{Remark 11.10.1.---}%
\footnote{\label{note:sga3-VIB-120}Editor's note: the first
part of the original \hyperref[rem:sga3-VIB-11.10.1]{Remark 11.10.1} has been incorporated
into \hyperref[lem:sga3-VIB-11.8]{11.8} and
\hyperref[prop:sga3-VIB-11.9]{11.9}
\textup{(with ``free'' replaced by ``projective'')}; the
counterexamples below correct its second part.}
In \hyperref[lem:sga3-VIB-11.8]{11.8}, it is not enough to
assume that \(C\) is a flat \(A\)-module, even when \(A\)
is a principal ideal domain. Indeed, we have the following
counterexamples, pointed out to us independently by
O.\ Gabber and J.-P.\ Serre.

\smallskip
\phantomsection\label{rem:sga3-VIB-11.10.1-a}%
\noindent\textup{(a)}
Let \((A,\mathfrak m)\) be a discrete valuation ring, let
\(K\) be its field of fractions, and let \(G\) be the
\(A\)-group obtained by ``extension by zero'' from the
\(K\)-group \((\mathbb Z/2\mathbb Z)_K\). The constant
group \((\mathbb Z/2\mathbb Z)_A\), and hence also its
subgroup \(G\), acts on the free \(A\)-module \(V\) with
basis \(v_1,v_2\) by interchanging \(v_1\) and \(v_2\).
Then \(Av_1\) is not a \(G\)-submodule of \(V\), but it is
the intersection of the \(G\)-submodules
\[
  Av_1+\mathfrak m^n v_2,\qquad n\in\mathbb N^*.
\]
Thus there is no smallest \(G\)-submodule of \(V\)
containing \(v_1\).

\smallskip
\phantomsection\label{rem:sga3-VIB-11.10.1-b}%
\noindent\textup{(b)}
Let \(A\) be an integral domain distinct from its field of
fractions \(K\), and let \(G\) be the flat affine
\(A\)-group corresponding to the Hopf algebra
\phantomsection\label{form:sga3-VIB-11.10.1-Hopf-example}%
\[
  \mathcal A(G)=\{P\in K[T]\mid P(0)\in A\}.
\]
Its comultiplication, counit, and antipode are respectively
defined by
\[
  \Delta(T)=T\otimes1+1\otimes T,\qquad
  \varepsilon(T)=0,\qquad
  \tau(T)=-T.
\]
\textup{(Thus \(G\otimes_AK=\mathbb G_{a,K}\).)}

Let \(V=Av_1\oplus Av_2\) be free over \(A\), and let
\(u\) be the endomorphism of \(V\) defined by
\[
  u(v_1)=v_2,\qquad u(v_2)=0,
\]
so that \(u^2=0\). Then \(V\) has the
\(\mathcal A(G)\)-comodule structure
\phantomsection\label{form:sga3-VIB-11.10.1-coaction}%
\[
  \mu(m)=m\otimes1+u(m)\otimes T.
\]
The \(G\)-submodules of \(V\) containing \(v_1\) are
precisely the \(A\)-submodules \(Av_1\oplus Iv_2\), where
\(I\) is a nonzero ideal of \(A\). Their intersection is
\(Av_1\), which is not a \(G\)-submodule. Hence there is no
smallest \(G\)-submodule of \(V\) containing \(v_1\).
Moreover,
\[
  C=A\oplus K\cdot T
\]
is a subcoalgebra of \(\mathcal A(G)\), flat over \(A\),
and the coaction
\[
  \mu:V\longrightarrow V\otimes\mathcal A(G)
\]
factors through \(V\otimes C\). Thus the ``very simple''
coalgebra \(C\) also gives a counterexample.

Finally, the preceding two examples are special cases of
the following construction. Let \(A\) be an integral
domain distinct from its field of fractions \(K\), and let
\(B\) be a Hopf \(A\)-algebra, free over \(A\). Write
\[
  \varepsilon:B\longrightarrow A,\qquad
  I=\ker(\varepsilon)
\]
for the augmentation and the augmentation ideal. Since
\(B=A\cdot1\oplus I\), one readily sees that
\phantomsection\label{form:sga3-VIB-11.10.1-B-prime}%
\[
  B'=\{b\in B\otimes_AK\mid \varepsilon(b)\in A\}
\]
is a Hopf subalgebra of \(B_K\). Suppose that \(V\) is a
\(B\)-comodule, free with basis
\((v_1,\ldots,v_n)\) as an \(A\)-module, and that
\(\mu_V(v_1)\neq v_1\otimes1\). Then \(Av_1\) is not a
subcomodule of \(V\), but it is the intersection of the
subcomodules
\[
  Av_1+J\cdot V
\]
as \(J\) ranges over the nonzero ideals of \(A\).%
\begingroup
\renewcommand{\thefootnote}{\fnsymbol{footnote}}%
\footnote[1]{\label{note:sga3-VIB-11.10.1-ideal-letter}
Translator's note: the French re-edition uses \(I\) both
for the fixed augmentation ideal above and for the ideals
ranging in this final intersection. We use \(J\) for the
latter to keep the two roles distinct.}%
\endgroup\ 
Thus there is no smallest subcomodule of \(V\) containing
\(v_1\).

\medskip
\phantomsection\label{prop:sga3-VIB-11.11}%
\noindent\textbf{Proposition 11.11.---}
\emph{Let \(G\) be an algebraic group over a field \(k\).
The following conditions are equivalent:}
\begin{enumerate}
\item[\textup{(i)}]\phantomsection\label{concl:sga3-VIB-11.11-i}
  \emph{\(G\) is affine.}
\item[\textup{(ii)}]\phantomsection\label{concl:sga3-VIB-11.11-ii}
  \emph{\(G\) is quasi-affine.}
\item[\textup{(iii)}]\phantomsection\label{concl:sga3-VIB-11.11-iii}
  \emph{\(G\) acts faithfully on a quasi-affine
  \(k\)-scheme \(X\).}
\item[\textup{(iv)}]\phantomsection\label{concl:sga3-VIB-11.11-iv}
  \emph{\(G\) acts linearly and faithfully on a
  \(k\)-vector space \textup{(not necessarily
  finite-dimensional)}.}
\item[\textup{(v)}]\phantomsection\label{concl:sga3-VIB-11.11-v}
  \emph{\(G\) is isomorphic to a closed subgroup of
  \(\operatorname{GL}(n)_k\).}
\end{enumerate}

\emph{Proof.}
The implication
\hyperref[concl:sga3-VIB-11.11-i]{\textup{(i)}}
\(\Rightarrow\)
\hyperref[concl:sga3-VIB-11.11-ii]{\textup{(ii)}} is
immediate, and
\hyperref[concl:sga3-VIB-11.11-ii]{\textup{(ii)}}
\(\Rightarrow\)
\hyperref[concl:sga3-VIB-11.11-iii]{\textup{(iii)}} because
\(G\) acts faithfully on itself by translations.

Suppose that \(G\) acts faithfully on the right on a
quasi-affine \(k\)-scheme \(X\).%
\footnote{\label{note:sga3-VIB-121}Editor's note: the
original has been expanded in what follows. We have also
chosen a right action of \(G\) on \(X\), so as to obtain a
left linear action of \(G\) on \(\mathcal O(X)\).}
Since \(X\) is quasi-affine, it is separated and
quasi-compact.%
Similarly, \(G\) is separated
\textup{(\hyperref[sga3:VIA:0.3]{VI\(_A\), 0.3})} and
quasi-compact \textup{(since it is of finite type over
\(k\))}. Hence, by
\hyperref[concl:sga3-VIB-11.1-c]{11.1\textup{(c)}}, there
are canonical isomorphisms
\phantomsection\label{form:sga3-VIB-11.11-functions-products}%
\[
  \mathcal O(X\times G)
    =\mathcal O(X)\otimes\mathcal O(G),
  \qquad
  \mathcal O(X\times G\times G)
    =\mathcal O(X)\otimes\mathcal O(G)\otimes\mathcal O(G).
\]
It follows that the morphism
\[
  \mu:\mathcal O(X)\longrightarrow
    \mathcal O(X)\otimes\mathcal O(G)
\]
induced by \(X\times G\to X\) gives
\(\mathcal O(X)\) the structure of a right
\(\mathcal O(G)\)-comodule; that is, \(G\) acts linearly on
the left on the \(k\)-algebra \(\mathcal O(X)\).
Consequently, \(G\) also acts on the right on the affine
envelope
\[
  X_{\mathrm{af}}=\operatorname{Spec}\mathcal O(X)
\]
of \(X\), and the canonical morphism
\(X\to X_{\mathrm{af}}\) is \(G\)-equivariant.

Since \(X\) is quasi-affine, the morphism
\(X\to X_{\mathrm{af}}\) is an open immersion
\textup{(EGA~II, 5.1.2)}. Hence, a fortiori, \(G\) acts
faithfully on \(X_{\mathrm{af}}\), and therefore the left
linear action of \(G\) on the \(k\)-algebra
\(\mathcal O(X)=\mathcal O(X_{\mathrm{af}})\) is faithful.
This proves
\hyperref[concl:sga3-VIB-11.11-iii]{\textup{(iii)}}
\(\Rightarrow\)
\hyperref[concl:sga3-VIB-11.11-iv]{\textup{(iv)}}.

Now suppose that \(G\) acts faithfully on a \(k\)-vector
space \(V\). By \hyperref[cor:sga3-VIB-11.10]{11.10},
\(V\) is the direct limit of finite-dimensional vector
subspaces \(V_i\) stable under \(G\). Let \(K_i\) be the
kernel of the induced action of \(G\) on \(V_i\); that is,
the kernel of \(G\to\operatorname{Aut}(V_i)\). Then \(K_i\)
is a closed subscheme of \(G\), and the faithfulness of the
action of \(G\) means that the intersection of the \(K_i\)
is the identity subgroup of \(G\). Since \(G\) is
Noetherian, one of the \(K_i\) is already the identity
subgroup. Thus \(G\to\operatorname{Aut}(V_i)\) is a
monomorphism, and hence, by
\hyperref[cor:sga3-VIB-1.4.2]{1.4.2}, a closed immersion.
This proves
\hyperref[concl:sga3-VIB-11.11-iv]{\textup{(iv)}}
\(\Rightarrow\)
\hyperref[concl:sga3-VIB-11.11-v]{\textup{(v)}}. Finally,
\hyperref[concl:sga3-VIB-11.11-v]{\textup{(v)}}
\(\Rightarrow\)
\hyperref[concl:sga3-VIB-11.11-i]{\textup{(i)}} is
immediate, and the proposition follows.

\medskip
\phantomsection\label{rem:sga3-VIB-11.11.1}%
\noindent\textbf{Remark 11.11.1.---}
Proposition \hyperref[prop:sga3-VIB-11.11]{11.11} admits
the following generalization. Let \(S\) be a regular,
locally Noetherian scheme of dimension at most \(1\), and
let \(G\) be a group scheme that is flat, quasi-compact, and
quasi-separated over \(S\). \textup{(In this case,
\(\mathcal A(G)\) is a torsion-free
\(\mathcal O_S\)-module, and hence is flat.)}

\begin{enumerate}
\item[\textup{(a)}]\phantomsection\label{concl:sga3-VIB-11.11.1-a}
  The following conditions are equivalent:%
  \begin{enumerate}
  \item[\textup{(i)}]\phantomsection\label{concl:sga3-VIB-11.11.1-a-i}
    \(G\) is affine over \(S\).
  \item[\textup{(ii)}]\phantomsection\label{concl:sga3-VIB-11.11.1-a-ii}
    \(G\) is quasi-affine over \(S\).
  \item[\textup{(iii)}]\phantomsection\label{concl:sga3-VIB-11.11.1-a-iii}
    \(G\) acts faithfully on a quasi-affine flat
    \(S\)-scheme.
  \item[\textup{(iv)}]\phantomsection\label{concl:sga3-VIB-11.11.1-a-iv}
    \(G\) acts linearly and faithfully on a flat
    quasi-coherent \(\mathcal O_S\)-module.
  \end{enumerate}

\item[\textup{(b)}]\phantomsection\label{concl:sga3-VIB-11.11.1-b}
  If, moreover, \(G\) is of finite type over \(S\) and
  \(S\) is Noetherian, these conditions imply that \(G\)
  is isomorphic to a closed subgroup of
  \(\operatorname{Aut}(\mathcal E)\), where
  \(\mathcal E\) is a finite-rank locally free
  \(\mathcal O_S\)-module.%
\end{enumerate}
